% !TEX TS-program = xelatex
% !TEX encoding = UTF-8 Unicode
% !Mode:: "TeX:UTF-8"

\documentclass{resume}
\usepackage{zh_CN-Adobefonts_external} % simplified Chinese Support using external fonts (./fonts/zh_CN-Adobe/)
\usepackage{linespacing_fix} % disable extra space before next section
\usepackage{cite}
\usepackage[colorlinks,linkcolor=blue]{hyperref}

\begin{document}
\pagenumbering{gobble} % suppress displaying page number


\name{张晓良}
\centerline{求职意向: 后端 / AI应用开发工程师 | 软件研发}
\vspace{1ex}
\contactInfo{
    {\href{mailto:zxl2024@qq.com}{zxl2024@qq.com}}}
    {\href{tel:+8615175951037}{(+86) 151-7595-1037}}
    {\url{https://github.com/HBUzxl} 
}

% \vspace{-1ex}
 
\section{\faGraduationCap\  教育背景}
\datedsubsection{\textbf{河北大学}~~ \ 硕士, 网络与信息安全, Top: 5\%}{2023.09 -- 2026.06}
\datedsubsection{\textbf{河北大学}~~ \ 学士, 计算机科学与技术, GPA: 4.1/5.0} {2019.09 -- 2023.06}

\vspace{1ex}


\section{\faUsers\ 实习经历}

\datedsubsection{\textbf{北京长亭科技有限公司~/~AI架构与探索组~/~实习AI应用开发工程师}}{2025.09 -- 2026.04}
% \vspace{-0.5ex}
\begin{onehalfspacing}
项目一：\href{https://monkeycode-ai.com/}{MonkeyCode} —— 企业级 AI 研发基础设施
\begin{itemize}
  \item 项目简介：该平台是集在线编程、代码生成、自动补全、代码分析与安全审计于一体的 AI 驱动研发平台，支持在线编程环境管理与任务执行，覆盖从需求设计到代码 Review 的全研发流程。负责分布式 Agent 的通信架构与编排，支撑大规模并发任务的高效运行。
  \item 核心通信架构设计：基于 gRPC 实现 Agent 与 Orchestrator 的双向流通信，设计了带指数退避算法的重连机制、心跳保活及状态上报，确保 Agent 在网络波动下能够实时同步任务队列与 AgentInfo
  \item 长连接回话与重放机制：实现 Agent 的终端会话（TTY）与文件管理系统 ，设计 ReconnectIO 会话重放机制，通过环形缓冲队列解决了web终端重连后的白屏断连问题
  \item 系统稳定性与性能调优：重构 Docker 构建上下文生成逻辑，利用 io.Pipe 实现流式 tar 打包，成功将环境创建时的内存峰值降低，解决了大规模并发下的 OOM 问题。
  \item 跨平台插件适配：主导 MonkeyCode 插件在 JetBrains 平台的适配，采用 RunVSAgent 方案复用 VSCode 插件代码，通过 JSON-RPC 模式抽离代码补全核心逻辑，实现跨平台体验高度一致。
  \item 外部 Agent 能力集成：基于 ACP 协议 将 Claude Code 等业内主流 Agent 接入平台，使用 Bun 封装 standalone 可执行文件，实现多系统、多架构的 Agent 运行环境保障。
\end{itemize}
  

项目二：\href{https://scan.monkeycode-ai.com/}{MonkeyScan} -- AI 驱动的代码安全扫描平台
\begin{itemize}
  \item 项目简介：企业级静态应用安全测试(SAST)平台，集成 CodeQL、Corax 等主流安全扫描引擎，为用户提供代码安全扫描、漏洞管理、合规审计等功能。负责系统的权限治理、三方授权集成、实名认证模块的开发。
  \item OAuth2/OIDC 认证：独立实现百智云单点登录集成，设计 OAuth State 防CSRF 攻击，通过Access Token / Refresh Token 及过期时间持久化，实现 Token 自动刷新与实名认证状态实时同步。
  \item 开发 GitHub App 授权模块：封装 ghinstallation 实现私有仓库的授权链路，设计 State 绑定用户身份的授权回调流程，支持私有仓库列表/分支查询，实现 Installation 绑定、JWT 签名校验与 Installation Access Token 自动生成，区分公开仓库与私有仓库的授权模式。
  \item 构建用户认证中间件链：实现 RequireCertification（实名认证检查）、RequireActivation（邀请码激活检查）、用户协议版本校验等多层拦截器，通过 Gin Context传递用户状态，实现了对敏感操作的细粒度权限控制与合规拦截。
  % \item 
  % \item 
  % \item 
  % \item 
  % \item 
\end{itemize}

  % 支持 GitHub/GitLab 仓库接入及 ZIP包上传，实现代码安全的全流程自动化管理。


  % -                                     
  % - 设计用户认证中间件链，实现实名认证检查、邀请码激活检查、用户协议版本校验
  % 2. 用户权限管理体系              
  % - 设计并实现用户角色权限系统，区分管理员/普通用户角色
  % - 开发用户管理后台 API（用户列表、角色更新、统计数据）
  % - 实现邀请码激活机制，支持用户注册流程管控
  % 3. 合规与风控功能
  % - 开发用户协议模块，支持协议版本管理与用户签署记录
  % - 集成百智云实名认证状态，实现实名用户权限控制（未实名用户限制敏感操作）
  % - 设计 Token 安全机制，支持 API Token 轮换与失效跳转
  % 4. 前端功能开发
  % - 开发登录页面、用户管理页面、实名认证提示弹窗等核心页面
  % - 实现任务列表提交人筛选器（管理员可见）
  % - 优化侧边栏布局、修复多项 UI 显示问题

  %   \item 基于 Go + Gin 构建后端服务，采用 Ent ORM 操作 PostgreSQL，使用 Wire 依赖注入、Viper 配置管理，通过 Redis实现会话存储与任务队列，整体服务通过 Docker Compose 编排部署。

  %   \item 设计邀请码激活与积分奖励机制，实现用户注册流程管控、邀请关系绑定、注册积分发放，通过事务保证邀请码消耗与积分变更的原子性，支持新用户首次登录引导与合规协议签署。
  %   \item 开发前端登录页面、用户管理页面、实名认证提示弹窗等核心组件，使用 React 19 + TypeScript + Vite 构建，实现任务列表管理员筛选器、侧边栏布局优化等功能。
\end{onehalfspacing}

\datedsubsection{\textbf{浙江远算科技有限公司~/~后端开发工程师}}{2025.06 -- 2025.08}
% \vspace{-0.5ex}
\begin{onehalfspacing}
\begin{itemize}
  \item 项目简介：3D模型智能配色生成工具，基于原有平台，开发了一套3D模型配色生成工具。用户可上传3D模型及相关参数，系统自动生成多样化配色方案，并支持任务管理与结果展示。
  \item 负责后端核心逻辑开发，包括任务创建、参数处理、图片解码与存储、异步任务调度等。
  \item 集成AI配色算法，自动生成符合场景和模板要求的配色方案。
  \item 实现任务状态管理与结果回调，优化用户体验。与前端及其他服务模块协作，保障系统稳定运行。 
\end{itemize}
\end{onehalfspacing}

\datedsubsection{\textbf{河北明智康华生物科技有限公司~/~研发工程师}}{2025.02 -- 2025.04}
% \vspace{-0.5ex}
\begin{onehalfspacing}
\begin{itemize}
  \item  设计并实现了一个基于前后端分离架构的远程会诊系统，提供用户认证、专家管理、病例登记、远程会诊、任务分配等功能，系统分为登录系统、管理系统、专家系统和分配系统四个模块
  \item 基于Go语言和Gin框架开发RESTful API接口，实现了用户鉴权、专家资源分配、病例管理等核心功能；前端使用Vue 3构建响应式用户界面
  \item 设计了基于JWT的认证机制，实现多角色(管理员、专家、医生)的权限控制和安全访问
  \item 使用CORS中间件解决跨域问题，保障前后端数据交互的安全性和稳定性
  % \item 开发了文件上传与管理模块，支持病例图像、诊断报告等医疗文档的存储和访问
  % \item 集成Swagger自动生成API文档，提升开发效率和团队协作
\end{itemize}
\end{onehalfspacing}





\vspace{-1ex}


\section{\faSitemap\ 项目经历}


\datedsubsection{\textbf{抖音电商(字节青训营)} }{2024.12 -- 2025.03}
\begin{onehalfspacing}
\begin{itemize}
  \item 设计并实现了一个基于微服务架构的完整电商平台，提供用户注册登录、商品浏览、购物车管理、订单结算和支付等功能,将系统拆分为用户、商品、购物车、结算、订单、支付等独立服务
  \item 设计了订单状态流转机制，通过RocketMQ实现订单超时自动取消、支付成功自动更新等功能,实现了购物车与订单的数据同步机制，实现订单、支付和库存等关键业务的数据一致性
  \item 设计并实现了基于Redis的库存预扣机制，解决高并发下的库存超卖问题
  \item 实现了基于EdDSA非对称加密的JWT认证机制，提高系统安全性
  \item 使用Snowflake算法实现分布式UUID生成，优化数据库查询性能
\end{itemize}
\end{onehalfspacing}
% \vspace{-1.5ex}

% \datedsubsection{\textbf{URL短链接生成器}} {2025.03 -- 2025.04}
% \begin{onehalfspacing}
    
% \begin{itemize}
%     \item 设计并实现了一个高性能的URL短链接服务，支持自定义短码生成、URL过期时间设置及过期链接自动清理功能，采用Go语言和Echo框架开发。
%     \item 集成Redis缓存层优化URL检索性能，显著减少数据库查询次数，提升系统响应速度
%     \item 实现了短码碰撞检测和自动重试机制，确保生成的短链接唯一可靠

% \end{itemize}
% \end{onehalfspacing}

% \vspace{-1ex}









\section{\faCogs\ 专业技能}
\begin{onehalfspacing}
\begin{itemize}

    \item 熟练掌握Go语言，理解slice、map等底层原理,了解GMP调度模型调度模型，垃圾回收等特性
    % \item 熟悉GMP调度模型、GC垃圾回收机制、内存分配与逃逸机制
    \item 熟悉Goroutine、Channel、Context，并能高效运行进行并发编程
    % \item 熟练使用Gin、Gorm、gRPC等常用框架，了解其内部实现原理
    \item 熟悉TCP/IP、UDP、HTTP/HTTPS等网络协议、理解TCP三次握手、四次挥手及IO复用机制
    \item 熟链表、栈、队列、树等常用数据结构，了解二分、动态规划、DFS、BFS、滑动窗口等常见算法
    \item 熟悉MySQL等关系型数据库，掌握存储引擎、索引结构、MVCC、锁机制、日志管理、查询优化及性能调优
    \item 熟悉Redis的数据结构、线程模型，了解持久化机制、缓存淘汰机制，了解缓存雪崩、穿透等问题
    \item 熟练使用Git进行代码管理，掌握GitHub/GitLab远程协作开发
    \item 熟悉Linux环境下编程，了解常用命令

\end{itemize}
\end{onehalfspacing}
% \vspace{-1ex}



\section{\faBook\ 学术成果}
\begin{onehalfspacing}
\begin{itemize}

\item \textbf{X. Zhang}, M. Wang, M. Wang, L. Wang, and H. Zhang, ``Enhanced RIS-assisted vehicular network with TDMA and Bayesian Compressive Sensing-based channel estimation,'' \emph{Computer Networks}, p. 111525, 2025. \textbf{(CCF-B)}~({\href{https://www.sciencedirect.com/science/article/abs/pii/S138912862500492X}{Link}})
\item \textbf{X. Zhang}, X. Zhai, M. Wang, and H. Zhang, ``Overlapping Community Detection Algorithm Based on Enhanced Label Propagation with Graph Neural Network Optimization,'' in \emph{Proceedings of the 21st International Conference on Intelligent Computing (ICIC 2025)}, Ningbo, China, Jul. 2025, pp. 2346--2357. \textbf{(CCF-C)}~({\href{https://scholar.google.com.sg/scholar?q=%E2%80%9COverlapping+Community+Detection+Algorithm+Based+on+Enhanced+Label+Propagation+with+Graph+Neural+Network+Optimization,&hl=zh-CN&as_sdt=0&as_vis=1&oi=scholart}{Link}})
\end{itemize}


\end{onehalfspacing}



\section{\faTrophy\ 奖项荣誉}
\begin{onehalfspacing}
\datedline{\textit{一等奖}~河北大学优秀学生奖学金}{2025.06}
\datedline{\textit{二等奖}~第六届高校研究生网络与信息安全技术大赛}{2023.09}
\datedline{\textit{一等奖}~全国大学生数学建模竞赛河北赛区}{2021.09}
\datedline{\textit{二等奖}~中国高校计算机大赛}{2022.04}
\datedline{\textit{}~CET6~/~CET4}{}
% \datedline{\textit{二等奖}~河北大学优秀学生奖学金}{2021.06}

\end{onehalfspacing}

\pagebreak

\end{document}
